\subsection{An inversion orbifold model for \texorpdfstring{$R7\_4\_2665$}{R7-4-2665}}
\label{sec:inversion-orbifold}

\begin{theorem}\label{thm:inversion-orbifold}
The ring $R7\_4\_2665$ has a unitary categorification.
\end{theorem}
\begin{proof}
Put $t=2+\sqrt5$. Start with Budinski's unitary near-group solution
$J_{16}^1$ for $C_{16}+16$, whose existence and exact certification are
given in \cite[Theorem~3.1.1, Table~3.1 and Section~3.3.2]{Budinski}.
Its data include
\[
 \chi(g,h)=\exp(2\pi i\,3gh/16),\qquad
 a(g)=\exp(-\pi i\,3g^2/16).
\]
Let $\mathcal C$ denote this category and $H=4C_{16}$. The central lifts
of its invertibles make $A=\bigoplus_{h\in H}h$ a commutative separable
algebra in $Z(\mathcal C)$: $H=H^\perp$ and $\chi|_{H\times H}=1$.
The category $\mathcal D=\mathcal C_A$ is unitary. Since $a(4)=-1$,
the orbifold rules in \cite[Theorem~12.5]{IzumiNear} are
\begin{equation}\label{eq:inversion-neutral}
 r^4=1,\qquad \rho r=r^{-1}\rho,\qquad
 \rho^2=r^2+(1+r+r^2+r^3)\rho.
\end{equation}
This is the already established rank-eight model $R8\_1\_12$.
Its canonical dual $\widehat H=C_4$ action has a generator $T$ fixing
$r$ and cycling the four dimension-$t$ simples.

\emph{The Morita symmetry and its central permutation.}
The regular $C_{16}$ algebra gives a self-Morita equivalence
$\mathcal M$ of $\mathcal C$ by \cite[Theorem~11.3]{IzumiNear}.
In the fixed-point model, with inclusion $\kappa:M^\alpha\to M$,
\[
 \alpha'_g=\operatorname{Ad}U(-g)|_{M^\alpha},\qquad
 \alpha_h\kappa=\kappa,\qquad
 \alpha_h(U(g))=\chi(h,g)U(g).
\]
The module has one fixed simple $a_0=\kappa$ of normalized dimension
four and sixteen simples $b_\lambda$ of dimension $t$ in a free orbit.
Its fixed simple has internal end the regular $C_{16}$ algebra, whose
structure is unique because $H^2(C_{16},\mathbb C^\times)=0$.
The opposite module has the same property. The argument of
\cite[Lemma~4.7 and Remark~4.8]{GJSCyclic} therefore makes
$\mathcal M\boxtimes_{\mathcal C}\mathcal M$ a regular bimodule with
a tensor autoequivalence $G_0$ on the right; uniqueness of this algebra
is the only part of that argument needed here.

The central permutation of $\mathcal M$ is inversion. Indeed,
$\operatorname{Ad}U(h)\rho=\rho\alpha_h$ gives
\[
 \operatorname{Ad}\rho(U(g))(U(h))=\chi(h,g)^{-1}U(h).
\]
On $p_\lambda=16^{-1}\sum_h\lambda(h)^{-1}U(h)$ in
$\operatorname{End}(\rho\kappa)$, left $\alpha_g$ and transported
right $\alpha'_g$ shift the character by $\chi_g^{-1}$ and
$\chi_g$, respectively. Every invertible has a unique central lift
by \cite[Theorem~13.4]{IzumiNear} and triviality of the universal
grading. The free orbit identifies its image as $g\mapsto-g$.
Thus the square fixes all central-invertible labels.
The faithful action on these labels in
\cite[Corollary~13.5]{IzumiNear} implies $G_0\simeq\mathrm{Id}$.
Inversion also preserves $H$ and its regular central algebra $A$.

\emph{Descent to a coherent involution.}
We retain the algebra identification needed for central base change.
An invertible bimodule whose center equivalence preserves $A$ descends,
after choosing that identification, to an invertible
$\mathcal C_A$--$\mathcal C_A$ bimodule of left $A$-module objects.
The identification balances the right action. This construction commutes
with relative tensor products: both sides classify balanced functors
with a compatible $A$-action, centrality moves $A$ across the balancing
variable, and separability splits the coequalizers. This comparison
respects associators and inverse bimodules.
Two algebra identifications differ by
$\operatorname{Aut}_{\rm alg}(A)=\widehat H$, and therefore by a
power of $T$ after descent. The descended bimodule $\mathcal L$ satisfies
\[
 \mathcal L\boxtimes_{\mathcal D}\mathcal L
 \simeq\mathcal D_{T^j}
\]
for some $j$ modulo four.

The object $a_0$ has four $A$-module structures, each of normalized
dimension one, and the free sixteen-element orbit gives four induced
simples of dimension $t$. Left invertibles act freely on the former.
The equation $\rho a_0=\sum_\lambda b_\lambda$ shows that the four
summands of $Q(\rho)$ take any chosen dimension-one module to the four
dimension-$t$ modules once each. Hence $\mathcal L$ is the regular
left $\mathcal D$-module, and its right action is an actual tensor
autoequivalence $F$.

Write its left $A$-action on $\kappa$ as scalars
$\lambda_h:\alpha_h\kappa\to\kappa$. Left $\alpha_g$ multiplies
them by $\chi(h,g)$; right $\alpha'_g$, transported by $U(g)$, gives
\[
 U(g)\lambda_h\alpha_h(U(g)^*)=\chi(h,g)^{-1}\lambda_h.
\]
This calculation on left $A$-modules is unaffected by the choice
balancing the right algebra. Therefore $F(r)=r^{-1}$. Adjusting that
choice by a character, and thus composing with $T$, gives
\begin{equation}\label{eq:inversion-F}
 F(r)=r^{-1},\qquad F(\rho)=r^{-1}\rho.
\end{equation}
At the tensor-functor level the regular-module identification retains
an inner ambiguity:
$F^2\simeq\operatorname{Inn}(r^\ell)T^j$.

Central monodromy removes it. The crossing of a central invertible
$p$ with an $A$-module has defect on the summand $h$ equal to
$\chi(p,h)^2$. Checking $A$-linearity of that crossing gives a
monoidal isomorphism between $\operatorname{Inn}(Q(p))$ and the dual
action of this character. The defect for $p=1$ is the order-two
character of $H$, whereas that for $p=2$ is trivial. Consequently
\[
 T^2\simeq\operatorname{Inn}(r),\qquad
 \operatorname{Inn}(r^2)\simeq\mathrm{Id}.
\]
By \eqref{eq:inversion-F}, $F^2$ fixes simple labels. Its action on
the noninvertibles gives $2\ell+j=0\pmod4$. Thus $j$ is even and
$F^2\simeq\operatorname{Inn}(r^{\ell+j/2})\simeq\mathrm{Id}$.
The universal grading of $\mathcal D$ is trivial by
\eqref{eq:inversion-neutral}, so
$\operatorname{Aut}_\otimes(\mathrm{Id}_{\mathcal D})=1$.
The two coherence composites for $F^3\to F$ therefore agree.
We obtain a coherent unitary $C_2$ action, not only an involution of
the based ring.

\emph{The reflection dual and its pointed group.}
Form $\mathcal E=\mathcal D\rtimes_F C_2$, with odd invertible $s$.
Its rules include
\[
 r^4=s^2=1,\quad srs=r^{-1},\quad
 \rho r=r^{-1}\rho,\quad \rho s=rs\rho,
\]
together with \eqref{eq:inversion-neutral}. Its eight invertibles form
$D_8$, and the other eight simples are a free orbit of dimension $t$.
The algebra $B=1+s$ exists from the coherent crossed product.
Let $\mathcal K={}_B\mathcal E_B$ and $L=\{1,s\}$.
On integral support, the double cosets $L$ and $r^2L$ each give two
invertibles; $LrL$ gives a unique dimension-two simple $X$.
The two irrational double orbits are
\[
 U\rho,\quad r^2U\rho,\qquad U=\{1,s,rs,r^{-1}\}.
\]
Their stabilizers are trivial, so each gives one simple, denoted $Y,Z$,
of dimension $2t$. Since $\rho^*=r^2\rho$, their duals interchange.

The group of four invertibles requires the associator, not just the
$D_8$ multiplication table. In the crossed-product factor model of
\cite[Section~12.1]{IzumiNear},
\[
 \alpha_1^4=\operatorname{Ad}\lambda_4,\qquad
 \alpha_1(\lambda_4)=\chi(1,4)\lambda_4=-i\lambda_4.
\]
The cyclic kernel obstruction is thus a primitive order-four class
in $H^3(C_4,\mathbb C^\times)$. The integral reflection dual is
$\mathcal K_{\rm int}=\mathcal C(D_8,L,\Omega,\psi)$.
The group-theoretical exponent formula
\cite[Theorem~9.2]{NgSchauenburgExponent} gives
\[
 \operatorname{FSexp}(\mathcal K_{\rm int})
 =\mathop{\rm lcm}_{C\le D_8\text{ cyclic}}
   |C|\,\operatorname{ord}(\Omega|_C).
\]
The subgroup $\langle r\rangle$ forces divisibility by $16$.
On the other hand $\mathcal K_{\rm int}$ is a Tambara--Yamagami
category, since it has four invertibles and a unique dimension-two
simple. Every such category on $V_4$ has Frobenius--Schur exponent
four or eight by \cite[Proposition~5.3]{LiuNgTotal}. The invertibles
of $\mathcal K$ therefore form $C_4=\langle g\rangle$.
Here $1,g^2$ lie above $L$ and fix $Y,Z$, while $g,g^3$ lie above
$r^2L$ and interchange them.

\emph{Complete fusion rules.}
The free induced bimodules are $X=BrB$, $Y=B\rho B$,
$Z=Br^2\rho B$. The genuine relative-tensor identity
\[
 (BuB)\otimes_B(BvB)\simeq BuBvB
\]
computes their products using the displayed rules of $\mathcal E$.
It gives
\begin{align*}
 X^2&=1+g+g^2+g^3,&
 XY=YX=XZ=ZX&=Y+Z,\\
 Y^2=Z^2&=g+g^3+X+4Y+4Z,&
 YZ=ZY&=1+g^2+X+4Y+4Z.
\end{align*}
For example the integral support of $Y^2$ is $2r^2L+LrL$;
Frobenius reciprocity gives $g,g^3$ once each, and the last part is
$X$. Every irrational constituent occurs four times. The other three
irrational products follow from the same identity. These rules,
$gX=X$, and $gY=Z$ determine all coefficients.
The independently generated model matches the full census tensor
under the permutation
\[
 (1,g,g^2,g^3,X,Y,Z)\longmapsto(0,1,6,2,5,3,4).
\]
All constructions preserve unitarity. This proves the theorem.
\end{proof}

The exact replay checks all $343$ source coefficients, all fusion axioms
of the rank-seven and rank-sixteen tensors, the neutral $R8\_1\_12$
identification, and the eight relative-induction products used above.
It does not infer the coherent action or the associator from those
finite checks; those are supplied by the preceding proof.
The same inversion-descent framework gives a second model with a
different parent group and cyclic associator.

\begin{theorem}\label{thm:inversion-v4}
The ring $R7\_4\_35$ has a unitary categorification.
\end{theorem}
\begin{proof}
Use Yu--Zhang's unitary $C_4\times C_4+16$ solution
\cite[Proposition~3.1 and Equation~(13)]{YuZhang}, with
\[
 \chi((x,y),(u,v))=i^{xu-yv},\qquad
 a(x,y)=\exp(2\pi i\,3(x^2-y^2)/8),\qquad c=1.
\]
The subgroup $H=\langle(1,1)\rangle$ is cyclic of order four,
$H=H^\perp$, and $a|_H=1$. Its quotient is generated by the class
$r$ of $(1,0)$. Central de-equivariantization gives
\begin{equation}\label{eq:inversion-neutral-selfdual}
 r^4=1,\qquad \rho r=r^{-1}\rho,\qquad
 \rho^2=1+(1+r+r^2+r^3)\rho.
\end{equation}

We justify the parent Morita square despite
$H^2(C_4\times C_4,\mathbb C^\times)\ne0$.
In the fixed-point model of Theorem~\ref{thm:inversion-orbifold},
the left internal-end algebra of $\kappa$ is untwisted. On the right,
choose $u_g=U(g):\kappa\alpha'_g\to\kappa$.
Right whiskering by $\alpha'_h$ precomposes the homomorphism and leaves
the implementing intertwiner unchanged. Hence the right stabilizer
cocycle is trivial as well:
\[
 u_h\circ(u_g\text{ right-whiskered by }\alpha'_h)
 =U(h)U(g)=U(g+h)=u_{g+h}.
\]
Here $U$ is the genuine representation in Izumi's normalization.
No application of $\alpha'_h$ to $U(g)$ is involved; $U(g)$ need not
belong to its domain $M^\alpha$.
Thus the original and inverse left modules have the same untwisted
regular internal-end algebra. Their square is regular up to a tensor
autoequivalence. The central permutation is again inversion, and its
square fixes all invertible labels. The faithful action in
\cite[Corollary~13.5]{IzumiNear} makes that autoequivalence trivial.

The central base-change and regularity arguments from
Theorem~\ref{thm:inversion-orbifold} now apply. The descended symmetry
can again be chosen with $F(r)=r^{-1}$ and $F(\rho)=r^{-1}\rho$.
Indeed $\chi((1,0),(j,j))=i^j$, so the left and right character
shifts are inverse and primitive. Their monodromy gives
$T^2\simeq\operatorname{Inn}(r)$, while $(2,0)$ gives a central
lift of $r^2$. These identities remove the same inner ambiguity in
$F^2$, and triviality of the universal grading from
\eqref{eq:inversion-neutral-selfdual} supplies coherence.
The unitary crossed product has eight invertibles $D_8$ and
$\rho s=rs\rho$, now with invertible term $1$ in $\rho^2$.

The reflection algebra $B=1+s$ again has a dual with four invertibles,
one dimension-two simple $X$, and two dimension-$2t$ simples $Y,Z$.
Here $\rho^*=\rho$, so both irrational double orbits are self-dual.
The pointed cyclic associator is trivial: $(1,0)$ is an actual
order-four complement to $H$, and its extended automorphism has fourth
power the identity, including on each $\lambda_{(j,j)}$ since
$(i^j)^4=1$.
Thus the associator $\Omega$ on the pointed $D_8$ part restricts
trivially to its rotation $C_4$. Every other cyclic subgroup has
order two, with a third-cohomology class of order at most two.
The formula in \cite[Theorem~9.2]{NgSchauenburgExponent} consequently
gives
\[
 \operatorname{FSexp}(\mathcal K_{\rm int})\mid4.
\]
The integral Tambara--Yamagami sector cannot have pointed group $C_4$.
For that group, \cite[Theorem~3.4 and its specialization]{ShimizuTY}
gives
\[
 \nu_4(X)=\frac12\sum_{x\in C_4}\chi(x,x)^{-1}=1+i
 \quad\text{or}\quad1-i,
\]
because nondegeneracy forces $\chi(1,1)=\pm i$. This differs from
$\FPdim(X)=2$, independently of the sign of the TY parameter,
contradicting the exponent bound. The pointed group is therefore
$V_4=\{1,g,h,gh\}$, with $h$ fixing $Y,Z$ and $g$ interchanging them.

The same free-induction identity yields the complete rules
\begin{align*}
 X^2&=1+g+h+gh,&XY=YX=XZ=ZX&=Y+Z,\\
 Y^2=Z^2&=1+h+X+4Y+4Z,&YZ=ZY&=g+gh+X+4Y+4Z.
\end{align*}
Together with $g^2=h^2=1$, $gX=hX=X$, $gY=Z$, and $hY=Y$,
these match all $343$ census coefficients under
\[
 (1,g,h,gh,X,Y,Z)\longmapsto(0,1,5,2,6,3,4).
\]
The separate source-independent replay verifies both tensor ranks,
all eight induction identities, and the fourth-indicator values.
The categorical operations are unitary, proving the theorem.
\end{proof}

The two constructions settle $R7\_4\_35$ and $R7\_4\_2665$ by one
descent method. They do not decide the other two sources
$R7\_4\_2053$ and $R7\_4\_2603$ in
Proposition~\ref{prop:four-dihedral}.
